% Shared four-copy circuit. Stages 1--7 derive the attack; 8--9 reduce it.
% Active columns use the same comfortable pitch, and each circuit is centered.
\newcommand{\distillationanalysisfigure}[1]{%
  \begin{tikzpicture}[
      wire/.style={thick,draw=CVQCInk},
      gate/.style={draw=CVQCInk,fill=white,thick,minimum width=1.0cm,
        minimum height=0.65cm,inner sep=2pt,font=\small,text=CVQCInk},
      pad/.style={gate,draw=CVQCTeal,fill=CVQCPale,minimum width=1.12cm},
      attack/.style={gate,draw=CVQCCoral,fill=CVQCCoralPale,
        text=CVQCCoral},
      input/.style={anchor=east,inner sep=0pt,font=\small,text=CVQCInk}
    ]
    \path[use as bounding box] (-6.6,-3.15) rectangle (6.6,3.05);

    \ifnum#1<5
      \def\distcolumns{7}
    \else
      \def\distcolumns{6}
    \fi
    \def\distpitch{1.55}
    \pgfmathsetmacro{\distfirst}{(2-(\distcolumns-1)*\distpitch)/2}
    \pgfmathsetmacro{\distlast}{\distfirst+(\distcolumns-1)*\distpitch}
    \pgfmathsetmacro{\distend}{\distlast+0.85}
    \pgfmathsetmacro{\distpt}{\distfirst}

    % Each build lays out only its active gate columns.
    \ifcase#1\relax
    \or % A_pt, OTP, A_ct, CNOTs, measurement, unmasking.
      \pgfmathsetmacro{\distpad}{\distfirst+\distpitch}
      \pgfmathsetmacro{\distct}{\distfirst+2*\distpitch}
      \pgfmathsetmacro{\distcnotone}{\distfirst+3*\distpitch}
      \pgfmathsetmacro{\distcnottwo}{\distfirst+4*\distpitch}
      \pgfmathsetmacro{\distmeter}{\distfirst+5*\distpitch}
      \pgfmathsetmacro{\distinverse}{\distfirst+6*\distpitch}
    \or % Move A_ct after the Clifford.
      \pgfmathsetmacro{\distpad}{\distfirst+\distpitch}
      \pgfmathsetmacro{\distcnotone}{\distfirst+2*\distpitch}
      \pgfmathsetmacro{\distcnottwo}{\distfirst+3*\distpitch}
      \pgfmathsetmacro{\distct}{\distfirst+4*\distpitch}
      \pgfmathsetmacro{\distmeter}{\distfirst+5*\distpitch}
      \pgfmathsetmacro{\distinverse}{\distfirst+6*\distpitch}
    \or % Commute the pad through the Clifford.
      \pgfmathsetmacro{\distcnotone}{\distfirst+\distpitch}
      \pgfmathsetmacro{\distcnottwo}{\distfirst+2*\distpitch}
      \pgfmathsetmacro{\distpad}{\distfirst+3*\distpitch}
      \pgfmathsetmacro{\distct}{\distfirst+4*\distpitch}
      \pgfmathsetmacro{\distmeter}{\distfirst+5*\distpitch}
      \pgfmathsetmacro{\distinverse}{\distfirst+6*\distpitch}
    \or % Move measured-copy unmasking before measurement.
      \pgfmathsetmacro{\distcnotone}{\distfirst+\distpitch}
      \pgfmathsetmacro{\distcnottwo}{\distfirst+2*\distpitch}
      \pgfmathsetmacro{\distpad}{\distfirst+3*\distpitch}
      \pgfmathsetmacro{\distct}{\distfirst+4*\distpitch}
      \pgfmathsetmacro{\distinverse}{\distfirst+5*\distpitch}
      \pgfmathsetmacro{\distmeter}{\distfirst+6*\distpitch}
    \or % Partial twirl, retaining the survivor's live pad.
      \pgfmathsetmacro{\distcnotone}{\distfirst+\distpitch}
      \pgfmathsetmacro{\distcnottwo}{\distfirst+2*\distpitch}
      \pgfmathsetmacro{\distpad}{\distfirst+3*\distpitch}
      \pgfmathsetmacro{\distct}{\distfirst+4*\distpitch}
      \pgfmathsetmacro{\distmeter}{\distfirst+5*\distpitch}
    \or % Follow a Pauli term of A_pt.
      \pgfmathsetmacro{\distcnotone}{\distfirst+\distpitch}
      \pgfmathsetmacro{\distcnottwo}{\distfirst+2*\distpitch}
      \pgfmathsetmacro{\distpad}{\distfirst+3*\distpitch}
      \pgfmathsetmacro{\distct}{\distfirst+4*\distpitch}
      \pgfmathsetmacro{\distmeter}{\distfirst+5*\distpitch}
    \or % Propagate the plaintext Pauli.
      \pgfmathsetmacro{\distcnotone}{\distfirst}
      \pgfmathsetmacro{\distcnottwo}{\distfirst+\distpitch}
      \pgfmathsetmacro{\distpropagated}{\distfirst+2*\distpitch}
      \pgfmathsetmacro{\distpad}{\distfirst+3*\distpitch}
      \pgfmathsetmacro{\distct}{\distfirst+4*\distpitch}
      \pgfmathsetmacro{\distmeter}{\distfirst+5*\distpitch}
    \or % Reduction: output Pauli is still beneath the live pad.
      \pgfmathsetmacro{\distcnotone}{\distfirst}
      \pgfmathsetmacro{\distcnottwo}{\distfirst+\distpitch}
      \pgfmathsetmacro{\distoutput}{\distfirst+2*\distpitch}
      \pgfmathsetmacro{\distpad}{\distfirst+3*\distpitch}
      \pgfmathsetmacro{\distct}{\distfirst+4*\distpitch}
      \pgfmathsetmacro{\distmeter}{\distfirst+5*\distpitch}
    \or % Commute that Pauli through the pad to the ciphertext level.
      \pgfmathsetmacro{\distcnotone}{\distfirst}
      \pgfmathsetmacro{\distcnottwo}{\distfirst+\distpitch}
      \pgfmathsetmacro{\distpad}{\distfirst+2*\distpitch}
      \pgfmathsetmacro{\distoutput}{\distfirst+3*\distpitch}
      \pgfmathsetmacro{\distct}{\distfirst+4*\distpitch}
      \pgfmathsetmacro{\distmeter}{\distfirst+5*\distpitch}
    \fi

    \foreach \y in {1.8,0.6,-0.6,-1.8} {
      \node[input] at ({\distfirst-1.05},\y)
        {$\mathsf{pTEnc}_{\pi}(\ket{0})$};
      \draw[wire] ({\distfirst-0.90},\y) -- ({\distmeter-0.30},\y);
      \draw[wire] ({\distfirst-0.80},{\y-0.13})
        -- ({\distfirst-0.70},{\y+0.13});
    }
    \draw[wire] ({\distmeter-0.30},1.8) -- (\distend,1.8);
    \foreach \y in {0.6,-0.6,-1.8} {
      \draw[wire,double,double distance=1.2pt]
        ({\distmeter+0.30},\y) -- (\distend,\y);
    }

    % CNOTs 1 -> 2 and 3 -> 4, followed by 3 -> 1.
    \draw[wire] (\distcnotone,1.8) -- (\distcnotone,0.6);
    \draw[wire] (\distcnotone,-0.6) -- (\distcnotone,-1.8);
    \draw[wire] (\distcnottwo,-0.6) -- (\distcnottwo,1.8);
    \foreach \x/\y in {\distcnotone/1.8,\distcnotone/-0.6,\distcnottwo/-0.6} {
      \fill[CVQCInk] (\x,\y) circle (0.055cm);
    }
    \foreach \x/\y in {\distcnotone/0.6,\distcnotone/-1.8,\distcnottwo/1.8} {
      \draw[wire,fill=white] (\x,\y) circle (0.12cm);
      \draw[wire] ({\x-0.12},\y) -- ({\x+0.12},\y);
      \draw[wire] (\x,{\y-0.12}) -- (\x,{\y+0.12});
    }

    \begin{scope}[xshift=\distmeter cm]
      \foreach \y in {0.6,-1.8} {
        \draw[wire,fill=white] (-0.30,{\y-0.24}) rectangle (0.30,{\y+0.24});
        \draw[wire] (-0.24,{\y-0.10})
          .. controls (-0.10,{\y+0.09}) and (0.10,{\y+0.09}) ..
          (0.24,{\y-0.10});
        \draw[wire] (-0.01,{\y-0.12}) -- (0.18,{\y+0.18});
      }
      \draw[wire,fill=white] (-0.30,-1.01) rectangle (0.30,-0.19);
      \draw[wire] (-0.24,-0.57)
        .. controls (-0.10,-0.38) and (0.10,-0.38) .. (0.24,-0.57);
      \draw[wire] (-0.01,-0.59) -- (0.18,-0.29);
      \node[inner sep=0pt,font=\small,text=CVQCInk] at (0,-0.82) {$X$};
    \end{scope}

    \ifnum#1<6
      \node[attack] at (\distpt,-0.6) {$A_{\mathrm{pt}}$};
    \fi
    \ifnum#1=6
      \node[attack] at (\distpt,-0.6) {$X^x Z^z$};
    \fi
    \ifnum#1<5
      \node[pad,minimum height=4.4cm] at (\distpad,0) {$\mathrm{OTP}$};
      \node[attack,minimum height=4.4cm] at (\distct,0)
        {$A_{\mathrm{ct}}$};
      \ifnum#1<4
        \node[pad,minimum height=3.2cm] at (\distinverse,-0.6)
          {$\mathrm{OTP}^{-1}$};
      \else
        \node[pad,minimum height=3.2cm] at (\distinverse,-0.6)
          {$\mathrm{OTP}^{\dagger}$};
      \fi
    \else
      % Only the measured pads cancel. The surviving pad stays live.
      \node[pad] at (\distpad,1.8) {$\mathrm{OTP}_1$};
      \node[attack] at (\distct,1.8) {$B_Q$};
      \ifnum#1<8
        \foreach \j/\y in {2/0.6,3/-0.6,4/-1.8} {
          \node[attack] at (\distct,\y) {$Q_{\j}$};
        }
      \else
        \node[attack] at (\distoutput,1.8) {$Q_4$};
      \fi
    \fi
    \ifnum#1=7
      \node[attack] at (\distpropagated,1.8) {$X^x$};
      \node[attack] at (\distpropagated,-0.6) {$X^x Z^z$};
      \node[attack] at (\distpropagated,-1.8) {$X^x$};
    \fi
  \end{tikzpicture}%
}

